%% preamble.tex — Налаштування документа для дипломної роботи (ДСТУ 3008:2015)

%% --- Мова та шрифти ---
\usepackage{fontspec}
\usepackage{polyglossia}
\setmainlanguage{ukrainian}
\setotherlanguage{english}
\setmainfont{Times New Roman}
\setsansfont{Arial}
\setmonofont[Script=Default,Scale=MatchLowercase]{Courier New}
\newfontfamily\cyrillicfonttt[Scale=MatchLowercase]{Courier New}

%% --- Геометрія сторінки ---
\usepackage[
  a4paper,
  left=3cm,
  right=1cm,
  top=2cm,
  bottom=2cm
]{geometry}

%% --- Міжрядковий інтервал ---
\usepackage{setspace}
\onehalfspacing

%% --- Абзацний відступ ---
\usepackage{indentfirst}
\setlength{\parindent}{1.25cm}

%% --- Вирівнювання тексту ---
\usepackage{microtype}

%% --- Терпимість до неперенесених монопросторних ідентифікаторів ---
\tolerance=2000                                  % default 200 — приймати «розтягнуті» рядки
\setlength{\emergencystretch}{3em}               % запас для горизонтального розтягу
\renewcommand{\_}{\textunderscore\allowbreak}    % дозвіл розриву після _ у texttt{}-ідентифікаторах

%% --- Запобігання вдовам/сиротам у верстці ---
\widowpenalty=10000        % немає 1-рядкового хвоста вгорі сторінки
\clubpenalty=10000         % немає 1-рядкового початку внизу сторінки
\displaywidowpenalty=10000 % те саме для абзаців, що закінчуються display-math
\widowpenalties=2 10000 10000  % також штрафуємо 2-рядкові «вдови»
\clubpenalties=2 10000 10000   % і 2-рядкові «сироти»

%% --- Заголовки ---
\usepackage{titlesec}

% Розділ (chapter): жирний, по центру, великі літери
\titleformat{\chapter}[block]
  {\normalfont\fontsize{14}{16}\bfseries\centering}
  {РОЗДІЛ \thechapter}
  {0pt}
  {\\\MakeUppercase}
\titlespacing{\chapter}{0pt}{0pt}{14pt}

% Підрозділ (section): жирний, ліворуч
\titleformat{\section}
  {\normalfont\fontsize{14}{16}\bfseries}
  {\thesection}
  {1em}
  {}
\titlespacing{\section}{0pt}{14pt}{7pt}

% Пункт (subsection): жирний, ліворуч
\titleformat{\subsection}
  {\normalfont\fontsize{14}{16}\bfseries}
  {\thesubsection}
  {1em}
  {}
\titlespacing{\subsection}{0pt}{14pt}{7pt}

%% --- Нумерація сторінок ---
\usepackage{fancyhdr}
\pagestyle{fancy}
\fancyhf{}
\fancyhead[R]{\thepage}
\renewcommand{\headrulewidth}{0pt}
\setlength{\headheight}{17pt}
\addtolength{\topmargin}{-2pt}

%% --- Таблиці ---
\usepackage{booktabs}
\usepackage{longtable}
\usepackage{array}
\usepackage{tabularx}
\usepackage{multirow}

%% --- Графіка ---
\usepackage{graphicx}
\usepackage{float}
\graphicspath{{./images/}}

%% --- Діаграми (TikZ) ---
\usepackage{tikz}
\usetikzlibrary{arrows.meta,positioning,calc,backgrounds,fit}

%% --- Підписи до рисунків і таблиць ---
\usepackage{caption}
\captionsetup{
  labelsep=endash,
  justification=centering,
  font={normalsize},
  labelfont=bf
}
\captionsetup[table]{position=top}
\captionsetup[figure]{position=bottom}

%% --- Списки ---
\usepackage{enumitem}
\setlist{noitemsep, topsep=0pt, partopsep=0pt}
\setlist[itemize]{leftmargin=1.25cm}
\setlist[enumerate]{leftmargin=1.25cm}

%% --- Математика ---
\usepackage{amsmath}
\usepackage{amssymb}

%% --- Лістинги коду ---
\usepackage{listings}
\usepackage{xcolor}

\definecolor{codegray}{rgb}{0.5,0.5,0.5}
\definecolor{codegreen}{rgb}{0,0.6,0}
\definecolor{codeblue}{rgb}{0,0,0.8}
\definecolor{backcolour}{rgb}{0.97,0.97,0.97}

\lstdefinestyle{mystyle}{
  backgroundcolor=\color{backcolour},
  commentstyle=\color{codegreen},
  keywordstyle=\color{codeblue},
  numberstyle=\tiny\color{codegray},
  stringstyle=\color{codegreen},
  basicstyle=\ttfamily\small,
  breakatwhitespace=false,
  breaklines=true,
  captionpos=b,
  keepspaces=true,
  numbers=left,
  numbersep=5pt,
  showspaces=false,
  showstringspaces=false,
  showtabs=false,
  tabsize=2,
  frame=single
}
\lstset{style=mystyle}

%% --- Гіперпосилання ---
\usepackage[
  hidelinks,
  unicode,
  pdftitle={Розробка VPN-застосунку з двофакторною автентифікацією},
  pdfauthor={},
  pdfsubject={Дипломна робота бакалавра},
  pdfkeywords={VPN, WireGuard, 2FA, TOTP, автентифікація}
]{hyperref}

%% --- Бібліографія ---
\usepackage[
  backend=biber,
  style=numeric,
  sorting=none,
  doi=true,
  url=true,
  maxbibnames=99
]{biblatex}
\addbibresource{references.bib}

%% --- Налаштування нумерації ---
\usepackage{chngcntr}
\counterwithin{figure}{chapter}
\counterwithin{table}{chapter}
\counterwithin{equation}{chapter}

%% --- Додаткові пакети ---
\usepackage{ifthen}
\usepackage{calc}
\usepackage{etoolbox}

%% --- Відступ після крапки в назвах ---
\usepackage{icomma}

%% --- Формат підписів для таблиць: «Таблиця N.M — Назва» ---
\DeclareCaptionLabelFormat{tablefmt}{Таблиця #2}
\captionsetup[table]{labelformat=tablefmt}

%% --- Формат підписів для рисунків: «Рисунок N.M — Назва» ---
\DeclareCaptionLabelFormat{figurefmt}{Рисунок #2}
\captionsetup[figure]{labelformat=figurefmt}

%% --- TODO-плейсхолдери ---
\newcommand{\TODO}[1]{\par\noindent\fbox{\textbf{<TODO>} #1}\par\medskip}
